\documentclass[11pt]{article}
\usepackage{amsmath,amssymb,amsthm,enumerate,nicefrac,fancyhdr,hyperref,graphicx,adjustbox}
\hypersetup{colorlinks=true,urlcolor=blue,citecolor=blue,linkcolor=blue}
\usepackage[left=2.6cm, right=2.6cm, top=1.5cm, includehead, includefoot]{geometry}
\usepackage[dvipsnames]{xcolor}
\usepackage[d]{esvect}

%% commands
%\include{commands.tex}

%% header
\pagestyle{fancy}
\fancyhead[L]{\bf\large CSC363 UTM \\ Assignment 2}
\fancyhead[R]{\bf\large Winter 2026 \\ Due Feb 24}
%\fancyfoot[C]{Page \thepage\ of 2}
\setlength{\headheight}{35pt}

\begin{document}

You are allowed to use AI, consult friends, TAs, but just do yourself a favor and understand what you are writing.  You are not allowed to post the assignment questions on any forums other than our Piazza.  

\vspace{0.2cm}

Please do not ask instructors to help with assignment questions. This is to ensure pedagogical equity and encourage independent problem solving.




\vspace{0.2cm}

\textbf{Every question on this assignment is worth 10 points.}


	\begin{enumerate}[{\bf Q1.}]



%Q1
\item  Prove that a set is semi-decidable iff it is $m$-reducible to $H=\{\left( e,x \right)\in\mathbb{N}^2 \colon M_e(x) \text{ halts}\}$. 

%Q2
\item Prove that every infinite semi-decidable set contains an infinite decidable subset.

%Q3
\item  Prove that if a set $A$ of natural numbers is semi-decidable and infinite, then there exists an injective computable function $f\colon\mathbb{N}\to\mathbb{N}$ such that $Range(f)=A$.

%Q4
\item Many standard G\"odel numbering schemes (such as those based on prime factorization or direct string encoding) are injective but not surjective -- leaving ``gaps" of natural numbers that do not correspond to valid Turing machine syntax. Prove that there exists a computable bijection between the set of G\"odel numbers of (valid syntax) Turing machines and the natural numbers.



%Q5
\item  Prove that the set $\mathbf{Tot} = \{e\in\mathbb{N} \colon M_e \text{ has domain }\mathbb{N}\}$ (i.e. the set of the indices of total computable functions) is not semi-decidable.

Your solutions must assume towards a contradiction that $\mathbf{Tot}$ is semi-decidable. Then, show that this assumption implies the existence of some computable $f$ such that $\{ {f(e)} \colon e \in\mathbb{N}\}$ is the set of indices of computable functions. This will imply that the function given by $g(x)=M_{f(x)}(x)+1$ is computable (explain why). Finally, find a contradiction.

\vspace{1cm}

%Task
For the next two questions, familiarize yourself with the following concepts: 

Partial Order Relations, Total Order Relations, Equivalence Relations, and Equivalence Classes of an equivalence relation. 

%Q6
\item Let  $\mathcal{P}(\mathbb{N})$ be the set of all subsets of $\mathbb{N}$ (the power set of $\mathbb{N}$). Prove that the relation $\leq_T$ on the set $\mathcal{P}(\mathbb{N})$ is a partial order but not a total order. 

%Q7
\item For $A,B$ sets of natural numbers, we write $A\equiv_T B$ if $A\leq_T B$ and $B\leq_T A$. Prove that $\equiv_T$ is an equivalence relation on $\mathcal{P}(\mathbb{N})$. What is the cardinality of the set of equivalence classes?



\newpage

%Q8
\item Prove that the Predecessor function is primitive recursive.


%Q9
\item
If $\mathcal{C}\subseteq \mathcal{P}(\mathbb{N})$, and $X\in\mathcal{C}$, we will say that $X$ is $\mathcal{C}$-complete if $\forall Y\in\mathcal{C}, Y\leq_m X.$

Prove that $\textbf{Fin}=\{ e  \in\mathbb{N}  \colon M_e \text{ halts on only finitely many elements}\}$ is $\Sigma^0_2$-complete. 

 

%Q10
\item  For $A,B$ sets of natural numbers, we will write $A\leq_p B$ if $A\leq_m B$ via a reduction function that is computable in polynomial time. More clearly, $A\leq_p B$ means that:

There exists $f\colon \mathbb{N}\to\mathbb{N}$ such that $f$ is computable in polynomial time, and $$\forall x\in\mathbb{N}, x\in A\Leftrightarrow f(x)\in B.$$

This new reducibility is known as \textbf{polynomial-time} reduction, which is important in studying complexity classes. 

Let $P = \{X\subseteq\mathbb{N}\colon X \text{ is decidable in polynomial time}\}$. 

\begin{enumerate}
    \item Let $U\in P$, and suppose that $X\leq_p U$. Prove that $X$ must be in $P$.
    \item Prove that, $\forall Y\in P, Y\leq_p \{1\}$.  
    \item $[$ True or False $]$ $\forall Y\in \Sigma^0_2, Y\leq_p \textbf{Fin}.$ Please justify.
    \item $[$ True or False $]$ $\forall Y\in \Sigma^0_1, Y\leq_p H$ (the halting set). Please justify.
\end{enumerate}


%Q11
\item \textbf{Bonus}

The last question part (b) shows that $\leq_p$ trivializes the concept of completeness for the class $P$, which is a very important class in Complexity Theory. In practice, $P$-completeness is based on  another type of reduction known as \textbf{Log-Space} reduction (denoted  $\leq_L$). We write $A\leq_L B$ if the reduction function is computable in $O(\log n)$ space (excluding the space for the read-only input and write-only output). Read about Log-space reduction, and the Circuit Value Problem (CVP). 

Prove that: CVP is $P$-complete (here $P$-completeness is with respect to $\leq_L$).

%Q12
\item \textbf{Bonus}

Prove that, for any given $e\in\textbf{Tot}$, there exists $n\in\mathbb{N}$ such that $M_{M_e(n)} = M_n.$
 

\end{enumerate}

\end{document}




    